\chapter*{Abstract}
\addcontentsline{toc}{chapter}{Abstract}
\vspace{-0.5cm}

From stars to galaxies, our Universe is observed to be in a magnetised state. In galaxies, for example, magnetic fields shape merger remnants, regulate star formation, and mediate the coupling of plasma motions on small and large scales. However, these fields are not expected to have existed in their current form in the hot, nearly uniform early Universe. Instead, weak seed fields were generated and subsequently amplified as galaxies assembled, where small-scale dynamos (SSDs) provided the most efficient route for field amplification. While SSDs have been explored extensively through analytic theories, laboratory experiments, and direct numerical simulations, the majority of our understanding is built on the assumption of incompressible turbulence. This limits direct application of SSD models to the supersonic, shock-dominated flow regime that characterises most of the interstellar medium.

In this thesis I use direct numerical simulations to extend SSD theory to compressible astrophysical regimes, and in the process develop novel statistical methods to extract robust, resolution-independent measurements from stochastic and intermittent flows. I apply these methods to two aspects of SSD evolution that are particularly relevant for galactic magnetism:
\begin{itemize}
    \item First, motivated by recent observations of young protogalaxies that reveal magnetic structures coherent on $\sim \mathrm{kpc}$ scales, I test whether compressibility modifies the scale on which magnetic fields become correlated. I show that shock-dominated turbulence shifts magnetic correlation scales from dissipation-regulated $\sim \mathrm{AU}$ scales to as large as $\sim \mathrm{pc}$ in some ISM phases. Compressibility, therefore, provides a pathway to generate substantially larger-scale (seed) fields, which must be further organised by (some yet, unidentified combination of) system-scale processes to reach the observed $\sim \mathrm{kpc}$ coherence.
    
    \item I then turn to the timescale over which SSDs saturate, and to the efficiency of nonlinear amplification once magnetic backreaction becomes important. Across both subsonic and supersonic regimes, I show that the SSD converts a nearly constant fraction of the turbulent kinetic energy flux into magnetic energy, with an efficiency consistent with that measured during the kinematic phase. This provides a direct empirical resolution to a long-standing debate regarding nonlinear SSD growth and suggests a wider universality of magnetic amplification.
\end{itemize}

Finally, as part of a longer-term campaign to enable SSD studies at higher resolution and larger dynamic range, I implement an ideal-MHD module in the new open-source, adaptive mesh refinement code \textsc{Quokka}, which was designed from the outset for next-generation GPU-accelerated supercomputers, and will enable the possibility of simulations with a much broader scale separation than is currently practical on CPU-based systems. This implementation couples a Godunov finite-volume scheme to a constrained-transport method that maintains $\nabla\cdot\mVector{b} = 0$ to numerical precision and incorporates stability and error-correction strategies specifically suited to highly compressible flows. This implementation provides a scalable foundation for future extensions that pursue more self-consistent, higher-resolution studies of SSD-driven magnetisation in galactic environments.
